% ============================================================
%  Розширення Related Work — три блоки
%  Файл: drafts/related-work-memory.tex
%  Дата: 2026-05-12
% ============================================================

% ----------------------------------------------------------
%  БЛОК 1.  Додаткові системи пам'яті (append до §2.1)
%  Вставити ПІСЛЯ рядка про Mem0 і ПЕРЕД абзацу «Обмеження:»
% ----------------------------------------------------------

A-MEM~\citep{xu2024amem} розвиває парадигму епізодичної пам'яті у напрямку структурованого
знання: кожне спостереження перетворюється на атомарну нотатку за принципом
Zettelkasten з явними зв'язками між нотатками, утворюючи граф знань поверх
потоку спостережень.
Це наближає одиницю пам'яті до структурованого запису, проте зв'язки
залишаються асоціативними, а не каузальними --- нотатка не фіксує
відхилені альтернативи та не прив'язується до конституційного принципу.

Letta~\citep{letta2025lettacode}, комерціалізований наступник MemGPT, вводить два механізми,
безпосередньо релевантні цій роботі.
По"=перше, \emph{контекстні репозиторії} (context repositories)~\citep{letta2026repositories} ---
персистентні версіоновані сховища структурованого контексту, до яких агент
звертається між сесіями; вони замінюють плоску архівну модель MemGPT
типізованими колекціями з підтримкою запитів.
По"=друге, \emph{контекстна конституція} (context constitution)~\citep{letta2026constitution} ---
декларативна специфікація того, до якого контексту агент має доступ і
з яким пріоритетом, аналогічна політиці витягування, вираженій як конфігурація,
а не як код.
Ці механізми наближають пам'ять агента до витягування за провенансом рішень:
контекстні репозиторії надають субстрат зберігання, а контекстна конституція ---
рудиментарну політику витягування.
Проте одиницею витягування в Letta залишається контекстний блок
(текстовий фрагмент з метаданими), а не оперативне рішення з провенансом,
альтернативами та конституційним обґрунтуванням.
Контекстна конституція визначає, \emph{з яких} колекцій запитувати,
а не \emph{який контекст рішення} витягти для поточної задачі.

Систематичний огляд механізмів пам'яті для LLM"=агентів подано
у~\citet{zhang2024survey}.

Усі розглянуті системи --- MemGPT, Generative Agents, A-MEM, Mem0, Letta ---
використовують як одиницю витягування діалоговий обмін, спостереження персонажа
або контекстний блок.
Контекстний блок ближчий до оперативного рішення, ніж діалоговий обмін,
проте не містить явних альтернатив, валідаторів та конституційного обґрунтування,
необхідних для провенансу рішень.
Ці системи ефективні для збереження ролі, відстеження преференцій та
епізодичного пригадування, але не для витягування одиниці, що визначає
довгострокову інженерну роботу: оперативного рішення з повним провенансом ---
яка альтернатива обрана, яка відхилена, чому, який валідатор забезпечує вибір
і який конституційний принцип його закріплює.


% ----------------------------------------------------------
%  БЛОК 2.  Синтез (новий абзац наприкінці §2)
%  Вставити ПІСЛЯ останнього підрозділу §2 (зараз §2.4)
% ----------------------------------------------------------

\medskip
\noindent\textbf{Синтез.}
Жодна з чотирьох ліній робіт --- діалогова пам'ять, витягування для програмних
агентів, реєстри архітектурних рішень (ADR), моделі з довгим контекстом ---
не розглядає провенанс рішення як першокласну одиницю пам'яті із семантичним
витягуванням.
Жодна не забезпечує примітиву повільного циклу оновлення (\emph{slow-loop refresh}),
що підтримує актуальність пам'яті про неактивні задачі на горизонтах у кілька
тижнів --- інтервалі, характерному для ротації операторів у штабних процесах.
Архітектура, описана в цій роботі, займає перетин цих чотирьох ліній:
вона запозичує структуру провенансу з ADR, семантичне витягування
з Code"=RAG, персистентність із систем діалогової пам'яті та селективне
завантаження контексту з RAG --- додаючи дворежимне витягування
(\emph{dual-mode retrieval}) як нову архітектурну примітиву, спеціально
спроектовану для забезпечення безперервності управління при зміні операторів.
